\section*{章末問題}
\begin{enumerate}[label=\textbf{\thechapter.\arabic*},leftmargin=*]
\item\label{問:面積と円環}
$y=x^2$ と $y=3x$ に囲まれた部分の面積を求めよ.また,この部分を $x$ 軸のまわりに回転させた立体の体積を求めよ.
\item\label{問:横帯の面積}
$x=y^2$ と $x=3-2y$ に囲まれた部分の面積を,\,$y$ を積分変数として求めよ.
\item\label{問:外半径の交代}
$0\leq x\leq2$,\,$-1\leq y\leq x$ の領域を $x$ 軸のまわりに回した立体の体積を求めよ.外半径が変わる位置を明示せよ.
\item\label{問:円筒と円板}
$0\leq x\leq1$,\,$x^2\leq y\leq1$ の領域を $y$ 軸のまわりに回した立体の体積を,円筒殻と円板の2通りで求めよ.
\item\label{問:弧長の平方完成}
$y=\frac13(x^2+2)^{\frac32}$,\,$0\leq x\leq1$ の曲線の長さを求めよ.
\item\label{問:道のりと平均速度}
$0\leq t\leq3$ における速度が $v(t)=3(t-1)(t-2)$,初期位置が $s(0)=0$ の運動について,終点の位置,道のり,平均速度,平均の速さを求めよ.
\item\label{問:平均値を取る点}
$f(x)=1+x^2$ の $[0,2]$ における平均値を求め,\,$f(c)$ がその平均値と等しくなる $c\in[0,2]$ を求めよ.
\item\label{問:三角形と両パップス}
頂点 $(0,0),(2,0),(2,1)$ の三角形を考える.面積重心の $y$ 座標を積分で求め,パップス＝ギュルダンの定理により,三角形を $x$ 軸のまわりに回した立体の体積を求めよ.また,境界の各線分の線重心を使って,できる立体の全表面積を求めよ.回転軸にある線分の寄与も説明せよ.
\end{enumerate}

\smallskip\noindent 解答は\pageref{解答:第9章}ページ.

